\oldpage[IV-2]{214}
\subsection{Anillos excelentes.}
\label{IV.2.ss7.8-fr}

\begin{env}[7.8.1]
\label{IV.2.s7.8.1-fr}
Los apartados anteriores de este parágrafo (así como (6.11), (6.12) y (6.13))
se han dedicado al estudio de problemas que aparecen con frecuencia al
utilizar anillos y preesquemas \emph{noetherianos}, y que pueden agruparse en los
tipos siguientes:

A) Para un anillo \emph{local} noetheriano $A$, ¿se «transmiten» las propiedades
de $A$ a su \emph{completado} $\widehat A$? Por ejemplo, si $A$ es reducido
(resp. un dominio de integridad, resp. un dominio de integridad integralmente
cerrado), ¿ocurre lo mismo con $\widehat A$? La mayoría de estas cuestiones
están relacionadas con las propiedades locales de las \emph{fibras formales}
de $A$ (recordemos que son las fibras del morfismo canónico
$\operatorname{Spec}(\widehat A)\to\operatorname{Spec}(A)$). Constituyen una
excepción las propiedades relacionadas con la noción de dimensión, por ejemplo
la de ser equidimensional; en ese caso, la \emph{condición de las cadenas} y
sus diversos refinamientos desempeñan el papel esencial.

B) Para un preesquema localmente noetheriano $X$ (en particular, para un esquema
afín $\operatorname{Spec}(A)$), ¿es \emph{abierto} el conjunto de los $x\in X$
en los que el anillo local $\mathcal O_x$ posee cierta propiedad (por ejemplo,
ser integralmente cerrado, de Cohen-Macaulay o regular)?

C) Para un anillo $A$ que es un \emph{dominio de integridad}, ¿es el cierre
entero de $A$ en una extensión finita de su cuerpo de fracciones un $A$-módulo
\emph{de tipo finito}? Naturalmente, esta cuestión puede formularse también
para los preesquemas noetherianos (\textbf{II}, 6.3).

Los problemas de tipo B) o C) pueden plantearse también para anillos
\emph{locales}, pero, en general, no basta que tengan una respuesta afirmativa
para cada anillo local $A_{\mathfrak p}$ de un anillo $A$ para que la tengan
para $A$ (cf. (6.13.6)).

Subrayemos, además, que al estudiar estos problemas nos hemos ocupado
sistemáticamente de determinar si una respuesta afirmativa es \emph{estable}
respecto de las dos operaciones más importantes del álgebra conmutativa:
la localización y el paso a una álgebra de tipo finito.

Los resultados obtenidos al estudiar estos problemas llevan a definir una
clase de anillos noetherianos cuyo comportamiento a este respecto es el mejor
posible:
\end{env}
